% !TEX root = ../main.tex
\begin{figure}[H]
  \centering
  \resizebox{\textwidth}{!}{%
    \begin{tikzpicture}[
      pantalla/.style={draw=diagNodoBorde, thick, rounded corners=4pt, fill=white},
      barra/.style={draw=diagGrupoBorde, fill=diagGrupo, rounded corners=2pt,
                    align=center, font=\footnotesize\bfseries, minimum height=0.8cm},
      usuario/.style={draw=diagNodoBorde, fill=diagNodo, rounded corners=4pt,
                      align=left, font=\scriptsize, inner sep=5pt},
      asistente/.style={draw=diagNodoBorde, fill=white, rounded corners=4pt,
                        align=left, font=\scriptsize, inner sep=5pt},
      fuentes/.style={draw=diagGrupoBorde, fill=diagGrupo, rounded corners=4pt,
                      align=left, font=\scriptsize, inner sep=5pt},
      entrada/.style={draw=diagNodoBorde, rounded corners=4pt, fill=white,
                      align=left, font=\scriptsize\itshape, inner sep=5pt},
      nota/.style={font=\scriptsize, text=diagGrupoBorde!60!black, align=left,
                   inner sep=2pt},
      apunte/.style={-{Latex[length=2mm]}, semithick, densely dashed,
                     draw=diagGrupoBorde!70!black},
      rotulo/.style={font=\small\bfseries, anchor=north},
      ]

      \node[pantalla, minimum width=8cm, minimum height=10.6cm] (p1) at (4,5.3) {};
      \node[barra, minimum width=7.4cm] at (4,9.9) {Asistente de titulaciones · EPSJ};
      \node[usuario, text width=4.6cm, anchor=north east] at (7.7,9.2)
      {¿En qué menciones se puede especializar el Grado en Ingeniería Eléctrica?};
      \node[asistente, text width=5.6cm, anchor=north west] at (0.3,7.7)
      {El Grado en Ingeniería Eléctrica tiene tres menciones:\\[2pt]
        · Generación Eléctrica con Energías Renovables\\
        · Instalaciones Eléctricas\\
        · Sistemas Eléctricos};
      \node[fuentes, text width=5.6cm, anchor=north west] at (0.3,4.5)
      {\textbf{Fuentes}\\[2pt]
        · Menciones del Grado en Ingeniería Eléctrica\\
        · Plan de estudios del Grado en Ingeniería Eléctrica};
      \node[entrada, text width=6.9cm, anchor=north west] at (0.3,1.6)
      {Escribe tu pregunta\ldots};
      \node[rotulo] at (4,-0.1) {Caso normal};

      \begin{scope}[shift={(9.6,0)}]
        \node[pantalla, minimum width=8cm, minimum height=10.6cm] (p2) at (4,5.3) {};
        \node[barra, minimum width=7.4cm] at (4,9.9) {Asistente de titulaciones · EPSJ};
        \node[usuario, text width=4.6cm, anchor=north east] at (7.7,9.2)
        {¿Se puede estudiar Enfermería aquí?};
        \node[asistente, text width=5.6cm, anchor=north west] (r2) at (0.3,7.7)
        {No he encontrado información sobre eso en la web de la Escuela
          Politécnica Superior de Jaén. Puedo ayudarte con las titulaciones que
          se imparten allí: sus asignaturas, qué se estudia en cada una y qué
          salidas profesionales tienen. ¿Sobre cuál te gustaría saber?};
        \node[nota, text width=5.4cm, anchor=north west] (n2) at (0.5,4.2)
        {Aquí \textbf{no} hay bloque de fuentes: no se ha recuperado ningún
          fragmento, así que no hay nada que citar. Su ausencia es la señal, y
          por eso no se sustituye por un bloque vacío.};
        \draw[apunte] (n2.north west) ++(0.2,0.15) -- ++(0,0.9);
        \node[entrada, text width=6.9cm, anchor=north west] at (0.3,1.6)
        {Escribe tu pregunta\ldots};
        \node[rotulo] at (4,-0.1) {Caso sin información suficiente};
      \end{scope}
    \end{tikzpicture}%
  }
  \caption[Guion gráfico de la pantalla de chat]{Guion
      gráfico de la pantalla de chat, en el caso normal y en el caso en que el
      sistema no encuentra información con la que responder.  Fuente: Elaboración propia.}
  \label{fig:storyboard_chat}
\end{figure}
